\section{Conclusions}

The horizon gate and the cross-covariance penalty split the persistent factor into a designated coordinate block, and across all six settings formed by three datasets and two loss families they jointly produce a separation that neither device reaches alone. The separation requires a latent target: switching only the prediction target to the raw signal halves the separation score, and what breaks down is precisely the inclusion of the persistent factor. On synthetic data we exceed label-free post-hoc rotation for both loss families, but on real data we fall short of Slow Feature Analysis, losing three of the four real settings, and the separation costs $0.111$ macro-F1 downstream on HAPT. The gate trades absolute performance for stability under domain shift.

Three measurements sharpen that comparison. Splitting SEP into its three factors shows that the shortfall against SFA is confined to the exclusion term --- the one term the objective cannot enforce --- while inclusion is a draw everywhere; that shortfall is not an artifact of linear probing, since it survives and in places widens under a non-linear reader which also collapses our own exclusion; and the coefficient that governs exclusion can be selected without task labels, agreeing with the label-based choice in all six settings.

\subsection{Limitations}

\textbf{The exclusion we demonstrate holds at the linear level only.} The objective enforces availability but not exclusion, and Section~\ref{sec:rotation} shows that this is not a formality: swapping the linear probe for a small MLP, with the same head given to every baseline, leaves inclusion and allocation intact and collapses exclusion alone, by $\numMlpExcl$ on average and $\numMlpExclWorst$ at worst. The transient factor is not absent from $\zslow$; it is linearly unreadable there. A structural repair does not help either --- a gate that also switches $\zslow$ off below $\tau$ is worse in $\numSymWorse$ of six settings. What we can claim is that the method fails where its own analysis said it would, and nowhere else.

\textbf{Most hyperparameters still cannot be chosen without labels.} The block width, the target width, the patch size and the horizon range have no derivation rule and were set by downstream performance, a known difficulty in this literature \citep{locatello2019}; the target width in particular governs how much of the transient factor is averaged away inside the target, and we never swept it. The cross-covariance coefficient is the exception, and a partial one: the label-free criterion $J$ agrees with the label-based choice in all six settings, but decisively only on synthetic data, where the margin to the runner-up is $\numJmarginSynthLo$--$\numJmarginSynthHi$ against $\numJmarginRealLo$--$\numJmarginRealHi$ on real data, and we did not store the seed-level spread that would settle the four real agreements. On a fourth dataset not reported here the criterion selected a different value from the label-based one; that analysis is unfinished --- one configuration axis was never tested --- and we report neither its results nor its disagreement as evidence.

\textbf{We did not demonstrate a practical benefit.} The perturbed-domain effect clears the error bars in only one of six settings, and even there absolute performance never overtakes the mechanism-free embedding: the gated block degrades less, which is not the same as being better.

\textbf{Our axis selects persistence, not relevance, and the case for it rests on data we built.} A persistent nuisance factor --- the low-frequency drift respiration and electrode motion produce in an electrocardiogram --- enters $\zslow$ by the same criterion, and which persistent factors matter remains a human judgment. The claim that a factor can be persistent yet absent from the slow part of the signal likewise rests on a single synthetic dataset: one counterexample does refute a universal proposition, but we designed this one ourselves, and neither PTB-XL nor HAPT exhibits the case.

\subsection{Future Work}

We compared only against label-free post-hoc rotation, so methods that place factor separation explicitly in the objective remain to be compared --- as does predictable feature analysis \citep{richthofer2013}, the post-hoc criterion closest to persistence and the one our argument in Section~\ref{sec:rotation} most owes a measurement. Validation is also needed on data whose persistent factor varies within a recording while staying constant within a window: PTB-XL fails the first property and HAPT the second, whereas sleep-stage classification has both. We have that fourth dataset in hand and the analysis is under way; it is the one referred to above, and it is unreported because its window-length axis has not been tested. Finally, reading the trajectory of the separated block as a state indicator rather than as classifier input (Appendix~\ref{app:traj}) would address the third limitation directly.
